% %%
\documentclass[lettersize,journal]{IEEEtran}
%\documentclass[lettersize,journal,draft]{IEEEtran}

\usepackage{amsmath,amsfonts}
\usepackage{algorithmic}
\usepackage{algorithm}
\usepackage{array}
\usepackage[caption=false,font=normalsize,labelfont=sf,textfont=sf]{subfig}
\usepackage{textcomp}
\usepackage{stfloats}
\usepackage{url}
\usepackage{verbatim}
\usepackage{graphicx}
\usepackage{cite}
\hyphenation{op-tical net-works semi-conduc-tor IEEE-Xplore}
\def\BibTeX{{\rm B\kern-.05em{\sc i\kern-.025em b}\kern-.08em
    T\kern-.1667em\lower.7ex\hbox{E}\kern-.125emX}}
\usepackage{balance}
\begin{document}
\title{Distributed Jacobian Estimation for Adaptive Multirobot Navigation to Critical Points in 2D Vector Fields}
\author{Christopher Waight, Christopher A. Kitts \IEEEmembership{Senior Member, IEEE}
\thanks{Manuscript received Month Day, Year; revised Month Day, Year.}
\thanks{C. Waight is a PhD Candidate in the Robotic Systems Laboratory, Department of Mechanical Engineering, Santa Clara University, Santa Clara, CA 95053 USA (e-mail: cwaight@scu.edu)}
\thanks{C. A. Kitts is Professor and Director of the Robotic Systems Laboratory, Department of Mechanical Engineering, Santa Clara University, Santa Clara, CA 95053 USA (e-mail: ckitts@scu.edu)}
}
\markboth{Journal of \LaTeX\ Class Files,~Vol.~XX, No.~X, Month~Year}%
{Waight, Kitts \MakeLowercase{\textit{et al.}: Adaptive Navigation of Multirobot Systems to Critical Points in 2D Vector Fields}}

%\IEEEpubid{0000--0000/00\$00.00~\copyright~2025 IEEE}

\maketitle

\begin{abstract}
    Multirobot systems operating within vector fields often need to locate and exploit key flow features such as sources, sinks, saddles, and vortex centers. Existing approaches for operating in vector fields either require prior knowledge of the environment or use local information to provide directional guidance. This paper presents a curvature aware adaptive navigation method using real time Jacobian estimates to enable precise critical point detection and control. We develop distributed control primitives requiring only three robots to estimate critical point locations and types through instantaneous local field measurements. The method operates without memory requirements or prior field knowledge and is scalable  beyond the minimum configuration. Simulation validation across six canonical vector fields confirms successful convergence and orbital control using three robot and four robot configurations. Hardware validation on the Decabot testbed demonstrates succcesful performance across 258 trials. In 236 convergence experiments, the system achieved a 100\% success rate with average position errors within 0.015 m of true critical points. An additional 22 orbital control experiments maintained commanded radii with standard deviations ranging from 0.01 m to 0.10 m. These results demonstrate that the approach also scales when additional robots are used.

\end{abstract}

\begin{IEEEkeywords}
Multirobot Systems, Adaptive Navigation, Vector Fields, Critical Points, Distributed Sampling, Cluster Control
\end{IEEEkeywords}

\section{Introduction}

\IEEEPARstart{M}{ultirobot} systems exploit the spatial distribution of robots to sample and compute local field properties, enabling monitoring and exploration of environments too hazardous for humans or detection animals [1--3]. Such systems have been deployed across diverse domains including flexible manufacturing, environmental monitoring, and search-and-rescue operations [4, 5]. Adaptive navigation approaches become necessary when operating in time-varying fields or environments with no prior knowledge [6--8]. The environments these systems navigate fall into two categories: scalar fields, such as temperature or chemical concentration, and vector fields, such as ocean currents, wind streams, and electromagnetic fields.

%\IEEEpubidadjcol

In scalar fields, every point in space is associated with a single value. Adaptive navigation techniques for multirobot systems enable robots to find extrema and saddle points, and to navigate along ridges, trenches, contour lines, and fronts [9, 32]. A variety of methods have been developed to locate these features, including gradient following, Hessian estimation, extended information consensus filters, constrained nonsmooth distributed optimization, and bio-inspired algorithms [11--14]. These techniques have been extensively validated for source seeking, plume tracking, and gradient following across ground rovers, quadrotors, and fixed-wing UAVs  [16--20].

Unlike scalar fields, vector fields possess both magnitude and direction at each point, producing vectors such as velocity in a wind stream. Vector fields can exhibit features absent in scalar fields such as separatrices, and critical points where the field magnitude of the flow field vanishes. Critical points can model phenomena such as plumes and vortices arising from tidal flows over submarine caves, and can be used to locate pollution sources[21,22]. Vector fields have also been used to construct navigation functions for obstacle avoidance and path planning, though these methods address different objectives than critical point localization [23]. Recent work has developed adaptive vector field guidance methods that compensate for unknown dynamics and environmental disturbances, though these approaches focus on path following rather than critical point estimation [24].

Current methods for navigating toward critical points typically extract scalar quantities such as magnitude from the vector field and then apply scalar field techniques [21, 25]. Passive approachess such as the vector-sum techniques instruct robots to follow the average perceived flow [25--27]. While effective for coarse directional guidance, these approaches cannot recover precise critical point locations, which is needed when active control is required and prior knowledge of the field is unavailable, such as maintaining precise standoff distances for data collection quality or safety [28--29]. Related work on target encirclement has achieved stable orbital formations using bearing-based control [15], but these methods require target position information to be available to at least some agents—an assumption that does not hold for field-induced critical points where no such prior knowledge exists.

For example, underwater robots monitoring an oceanic vortex or atmospheric drones tracking a developing cyclone must estimate the field center to function effectively. Without localizing this critical point, these systems cannot maintain stable orbits at specified radii or compensate for flow-induced drift. While they may be able to converge toward or diverge from a feature, they lack the metric information necessary to determine their precise standoff distance.

This paper introduces a distributed estimation strategy that reconstructs local field geometry from instantaneous vector measurements. Using a minimum of three robots, our approach estimates the local Jacobian matrix, revealing both critical point locations and their types through eigenvalue analysis. This enables precise navigation primitives including controlled approach trajectories and stable orbital control at fixed radii.

Our contributions are fourfold. First, we present a distributed algorithm that estimates equilibrium locations and types using only instantaneous measurements. Second, we show the approach scales beyond the minimum configuration to work with additional robots without algorithmic changes. Third, we develop control strategies that exploit these estimates for precise navigation to and around critical points across six vector field types. Fourth, we validate our approach through 258 hardware experiments demonstrating repeatable critical point estimation and orbit control.

The rest of this paper details our contributions. Section II develops the mathematical framework for distributed Jacobian estimation and presents control primitives for each field singularity type. Section III describes the control architecture used to verify these primitives. Section IV presents simulation results, while Section V provides hardware validation. Section VI analyzes performance relative to other primitives and discusses the impact, limitations, and future extensions of this work. Section VII concludes.



\section{Critical Point Detection and Control}

\subsection{Vector Field Environments}

A vector field assigns a vector $\mathbf{v}(\mathbf{p}) = (u(\mathbf{p}), v(\mathbf{p}))$ to each point $\mathbf{p} = (x, y)$ in a planar domain. In physical systems, this represents directional quantities such as fluid velocity, electromagnetic force, or wind patterns.

A critical point $\mathbf{p}^*$ occurs where the field vanishes:
\begin{equation}
\mathbf{v}(\mathbf{p}^*) = \mathbf{0}
\end{equation}

Near critical points, the field behaviour is characterized by the Jacobian matrix $\mathbf{J} = \nabla\mathbf{v}|_{\mathbf{p}}$, defined at all points $\mathbf{p}$ (not just $\mathbf{p}^*$). The eigenvalues of the Jacobian evaluated at the critical point $\mathbf{p}^*$ determine the topological type. Six canonical critical point types arise from distinct eigenvalue structures, as shown in Fig.~\ref{fig:allfield} and classified in Table~\ref{tab:field_summary}.



\begin{figure}[!t]
    \centering
    \includegraphics[width=0.37\textwidth]{six_vector_fields.png}
    \caption{Six linear vector fields types with critical points}
    \label{fig:allfield}
\end{figure}

\begin{table}[!t]
    \centering
    \caption{Classification of vector field critical points based on Jacobian eigenvalue structure ($\lambda$ = eigenvalue; complex eigenvalues $\alpha \pm i\omega$ where $\alpha$ is real part, $\omega$ is imaginary part).}
    \label{tab:field_summary}
    \begin{tabular}{lll}
    \hline
    \textbf{Field Type} & \textbf{Critical Point} & \textbf{Eigenvalues} \\
    \hline
    Vortex & Center & $\lambda = \pm i\omega$ \\
    Sinking Vortex & Stable Spiral & $\lambda = \alpha \pm i\omega$, $\alpha < 0$ \\
    Sink & Stable Node & $\lambda_1, \lambda_2 < 0$ \\
    Source & Unstable Node & $\lambda_1, \lambda_2 > 0$ \\
    Spewing Vortex & Unstable Spiral & $\lambda = \alpha \pm i\omega$, $\alpha > 0$ \\
    Saddle & Saddle Point & $\lambda_1 < 0 < \lambda_2$ \\
    \hline
    \end{tabular}
    \end{table}

Eigenvalue analysis can also identify scenarios such as repeat roots, or the determinant might be 0 ( Uniform flow fields, for instance, have constant velocity everywhere with $\mathbf{J} = \mathbf{0}$.) In these cases, no accessible critical point exists. These scenarios are detectable through our estimation framework but are not considered further in this work.


\subsection{Critical Point Detection and Identification}
\noindent

Given a cluster of three robots sampling a planar vector field with robot positions $\mathbf{p}_i = (x_i, y_i)$ and corresponding vector field measurements $(u_i, v_i)$ for $i \in \{1, 2, 3\}$, we can approximate the vector field components across the cluster's spatial extent using linear planar models:
\begin{equation}
U(x,y) = ax + by + c
\end{equation}
\begin{equation}
V(x,y) = dx + ey + f
\end{equation}

The coefficients are determined by solving the linear systems:
\begin{equation}
\mathbf{A}\boldsymbol{\theta}_u = \mathbf{u}, \quad \mathbf{A}\boldsymbol{\theta}_v = \mathbf{v}
\end{equation}
where:
\begin{equation}
\mathbf{A} = \begin{bmatrix} x_1 & y_1 & 1 \\ x_2 & y_2 & 1 \\ x_3 & y_3 & 1 \end{bmatrix}, \quad \boldsymbol{\theta}_u = \begin{bmatrix} a \\ b \\ c \end{bmatrix}, \quad \boldsymbol{\theta}_v = \begin{bmatrix} d \\ e \\ f \end{bmatrix}
\end{equation}
and $\mathbf{u} = [u_1, u_2, u_3]^T$, $\mathbf{v} = [v_1, v_2, v_3]^T$ are the collected field measurements.

The critical point $\mathbf{p}^* = (x^*, y^*)$ satisfies $U(x^*, y^*) = 0$ and $V(x^*, y^*) = 0$, which yields:
\begin{equation}
ax^* + by^* + c = 0, \quad dx^* + ey^* + f = 0
\end{equation}

Rewriting in matrix form:
\begin{equation}
\mathbf{J}\mathbf{p}^* = -\mathbf{h}
\end{equation}
where the Jacobian matrix and offset vector are:
\begin{equation}
\mathbf{J} = \begin{bmatrix} a & b \\ d & e \end{bmatrix}, \quad \mathbf{h} = \begin{bmatrix} c \\ f \end{bmatrix}
\end{equation}

Thus, provided $\det(\mathbf{J}) \neq 0$:
\begin{equation}
\mathbf{p}^* = -\mathbf{J}^{-1}\mathbf{h}
\end{equation}

The eigenvalues of the Jacobian can be used to identify the type of nearby critical point (Table 1).

\subsection{Scaling Beyond Minimal Configurations}

To incorporate additional robots beyond the minimum three, we form an overdetermined system using all measurements from all robots in the cluster.

With $n > 3$ robots, the system of equations in equation 3 becomes overdetermined. We form the extended measurement matrix:
\begin{equation}
\mathbf{A} = \begin{bmatrix} x_1 & y_1 & 1 \\ \vdots & \vdots & \vdots \\ x_n & y_n & 1 \end{bmatrix} \in \mathbb{R}^{n \times 3}
\end{equation}

Since $n > 3$, we employ least-squares minimization:
\begin{equation}
\boldsymbol{\theta}_u = \arg\min_{\boldsymbol{\theta}} \|\mathbf{A}\boldsymbol{\theta} - \mathbf{u}\|^2 = (\mathbf{A}^T\mathbf{A})^{-1}\mathbf{A}^T\mathbf{u}
\end{equation}

Similarly for $\boldsymbol{\theta}_v$. The critical point is then computed using the same Jacobian formulation as before. This approach considers all measurements from robotic sensors  and provides improved estimation through redundancy, with robustness increasing for uncorrelated measurement errors.

% \subsection{Critical Point Attraction Control Law}
% The average position of all robots in a cluster is known as $\mathbf{p}_c$, the cluster center. Given the current cluster center $\mathbf{p}_c = (x_c, y_c)$ and the estimated critical point $\mathbf{p}^*$, the continuous velocity command is:
% \begin{equation}
% \dot{\mathbf{p}}_c = k(\mathbf{p}^* - \mathbf{p}_c)
% \end{equation}
% where $k > 0$ is the proportional gain. This yields exponential convergence to the estimated critical point with time constant $\tau = 1/k$.


% \subsection{Orbital Navigation Strategy}

% Given the estimated critical point $\mathbf{p}^*$ from equation (8), we design a control law to maintain a circular orbit at a desired radius $r_d$. The vector from the cluster center to the critical point is:

% \begin{equation}
% \mathbf{r} = \mathbf{p}^* - \mathbf{p}_c
% \end{equation}

% with magnitude:
% \begin{equation}
% r = \|\mathbf{r}\| = \sqrt{(x^* - x_c)^2 + (y^* - y_c)^2}
% \end{equation}

% The unit radial vector is:
% \begin{equation}
% \hat{\mathbf{r}} = \frac{\mathbf{r}}{r}
% \end{equation}

% The clockwise perpendicular vector is obtained through a $-90^\circ$ rotation:
% \begin{equation}
% \hat{\mathbf{r}}_\perp = \begin{bmatrix} \cos(-\pi/2) & -\sin(-\pi/2) \\ \sin(-\pi/2) & \cos(-\pi/2) \end{bmatrix} \hat{\mathbf{r}} = \begin{bmatrix} 0 & 1 \\ -1 & 0 \end{bmatrix} \hat{\mathbf{r}}
% \end{equation}


% The orbital velocity command combines tangential motion with radial correction:

% \begin{equation}
% \dot{\mathbf{p}}_c = v_t \hat{\mathbf{r}}_\perp + v_r \hat{\mathbf{r}}
% \end{equation}

% where the tangential velocity maintains orbital motion:
% \begin{equation}
% v_t = k_t
% \end{equation}

% and the radial velocity corrects for radius error:
% \begin{equation}
% v_r = -k_r(r - r_d)
% \end{equation}


% Substituting equations (15) and (16) into (14):

% \begin{equation}
% \dot{\mathbf{p}}_c = k_t \hat{\mathbf{r}}_\perp - k_r(r - r_d)\hat{\mathbf{r}}
% \end{equation}

% where $k_t > 0$ sets the orbital speed and $k_r > 0$ determines the radius correction rate. This control law produces stable circular orbits around the estimated critical point, with clockwise rotation for $k_t > 0$ (counterclockwise for $k_t < 0$) and exponential convergence to the desired radius with time constant $\tau_r = 1/k_r$.


% \subsection{Stability Analysis}

% We analyze the stability of both control modes and their robustness to critical point estimation errors.

% \subsubsection{Critical Point Attraction Stability}

% For the attraction control law $\dot{\mathbf{p}}_c = k(\mathbf{p}^* - \mathbf{p}_c)$, define the error vector:
% \begin{equation}
% \mathbf{e} = \mathbf{p}_c - \mathbf{p}^*
% \end{equation}

% The error dynamics are:
% \begin{equation}
% \dot{\mathbf{e}} = \dot{\mathbf{p}}_c = -k\mathbf{e}
% \end{equation}

% Consider the Lyapunov function:
% \begin{equation}
% V_1 = \frac{1}{2}\mathbf{e}^T\mathbf{e} = \frac{1}{2}\|\mathbf{e}\|^2
% \end{equation}

% Its time derivative along system trajectories:
% \begin{equation}
% \dot{V}_1 = \mathbf{e}^T\dot{\mathbf{e}} = -k\mathbf{e}^T\mathbf{e} = -k\|\mathbf{e}\|^2 < 0 \quad \forall \mathbf{e} \neq \mathbf{0}
% \end{equation}

% Since $\dot{V}_1 < 0$ for all $\mathbf{e} \neq \mathbf{0}$ and $k > 0$, the system exhibits global exponential stability with convergence rate $k$.

% \subsubsection{Orbital Control Stability}

% For the orbital control law, we transform to polar coordinates $(r, \theta)$ centered at $\mathbf{p}^*$, assuming $r > 0$ (i.e., the cluster center does not coincide with the estimated critical point). The dynamics become:
% \begin{equation}
% \dot{r} = v_r = -k_r(r - r_d)
% \end{equation}
% \begin{equation}
% \dot{\theta} = \frac{v_t}{r} = \frac{k_t}{r}
% \end{equation}


% \paragraph{Radial Convergence.}
% Define the radial error $e_r = r - r_d$ and consider the Lyapunov function:
% \begin{equation}
% V_2 = \frac{1}{2}e_r^2
% \end{equation}

% Its derivative along trajectories:
% \begin{equation}
% \dot{V}_2 = e_r \dot{e}_r = e_r \dot{r} = -k_r e_r^2 < 0 \quad \forall\, e_r \neq 0
% \end{equation}

% This establishes exponential convergence $r(t) \to r_d$ with time constant $\tau_r = 1/k_r$.

% \paragraph{Angular Dynamics.}
% The angular velocity $\dot{\theta}(t) = k_t / r(t)$ varies during the radial transient. As $r(t) \to r_d$, the angular velocity converges to the constant $\omega_\infty = k_t / r_d$, yielding steady circular motion at the desired radius.

% \subsubsection{Robustness to Estimation Errors}

% The control laws demonstrate robustness to errors in the estimated critical point location. Let $\boldsymbol{\epsilon} = \mathbf{p}^*_{est} - \mathbf{p}^*_{true}$ denote the estimation error. For the attraction controller, the system converges exponentially to the estimated critical point $\mathbf{p}^*_{est}$ with rate $k$, even when $\mathbf{p}^*_{est} \neq \mathbf{p}^*_{true}$. Similarly, the orbital controller maintains stable circular motion around $\mathbf{p}^*_{est}$ at the desired radius $r_d$. Thus, estimation errors shift the equilibrium location but preserve the exponential convergence properties established above.

% \subsection{Estimation Accuracy and Sensitivity Analysis}

% This section analyzes how formation geometry, sensor noise, and field characteristics affect critical point estimation accuracy.

% \subsubsection{Estimation Accuracy}

% The plane-fitting algorithm assumes local linearity of the vector field. For a globally linear field, three robots at any finite separation yield exact critical point estimates, as the plane fit recovers the true field coefficients. Real fields, however, are only locally linear. The approximation quality thus depends on formation size relative to field curvature.

% Define the formation size $s$ as the maximum distance between any two robots (e.g., the side length for an equilateral triangle). Through Taylor expansion, the neglected quadratic terms in the field approximation have magnitude $\mathcal{O}(s^2\|\mathbf{H}\|)$, where $\mathbf{H}$ is the field's Hessian tensor. These neglected terms propagate through the critical point estimation $\mathbf{p}^* = -\mathbf{J}^{-1}\mathbf{h}$, where $\mathbf{J}$ and $\mathbf{h}$ are the Jacobian matrix and offset vector defined in equations (6)--(7). Note that while the Hessian $\mathbf{H}$ characterizes local field curvature and thus the approximation error magnitude, the Jacobian $\mathbf{J}$ determines how this error maps to the critical point estimate. The resulting estimation error scales as:
% \begin{equation}
% \|\mathbf{p}^*_{est} - \mathbf{p}^*_{true}\| \sim \mathcal{O}(s^2)
% \end{equation}

% This quadratic scaling implies that halving the formation size reduces estimation error by approximately 75\%. For our experimental formations with $s = 0.33$ m, the nonlinearity-induced errors remain below 0.01 m, as confirmed by our hardware results.

% \subsubsection{Sensitivity to Measurement Noise}

% Robot sensors inevitably introduce measurement noise into the field readings. When each robot measures the field with additive noise $\boldsymbol{\eta}_i$, the resulting measurements become $\tilde{\mathbf{f}}_i = \mathbf{f}_i + \boldsymbol{\eta}_i$, where $\mathbf{f}_i = (u_i, v_i)$ denotes the true field vector at robot $i$.

% This measurement noise affects the plane-fitting coefficients and subsequently the critical point estimate. The error amplification depends on two geometric factors:
% \begin{equation}
% \|\Delta\mathbf{p}^*\| \propto \kappa(\mathbf{A}) \cdot \kappa(\mathbf{J}) \cdot \|\boldsymbol{\eta}\|
% \end{equation}

% Here, $\kappa(\mathbf{A})$ is the condition number of the robot formation geometry, and $\kappa(\mathbf{J})$ is the condition number of the local field Jacobian. Well-conditioned systems (small $\kappa$) provide stable estimates, while ill-conditioned systems amplify noise significantly.

% \subsubsection{Formation Geometry and Numerical Conditioning}

% The formation geometry directly impacts estimation stability through the condition number $\kappa(\mathbf{A})$. Among all triangle configurations, the equilateral triangle minimizes $\kappa(\mathbf{A})$, making it the optimal choice for robust estimation.

% For robots in an equilateral triangle formation with side length $s$ centered at the origin:
% \begin{equation}
% \mathbf{A}_{eq} = \begin{bmatrix}
% 0 & s/\sqrt{3} & 1 \\
% -s/2 & -s/(2\sqrt{3}) & 1 \\
% s/2 & -s/(2\sqrt{3}) & 1
% \end{bmatrix}, \quad \kappa(\mathbf{A}_{eq}) = \sqrt{6} \approx 2.45
% \end{equation}

% This optimal conditioning explains why equilateral formations consistently outperform other triangle geometries in our experiments. For four-robot configurations, a square formation with side length $s$ achieves $\kappa(\mathbf{A}_{square}) \approx 2.16$, a 12\% improvement over the triangle. This improved conditioning arises from better directional isotropy: the square's four corners provide more uniform spatial sampling with 90-degree symmetry. However, the four-robot overdetermined system requires a pseudoinverse solution rather than direct matrix inversion, introducing additional noise propagation paths that can degrade precision despite improved geometric conditioning.

% In contrast, degenerate configurations where robots approach collinearity yield $\kappa(\mathbf{A}) \to \infty$, causing the estimation to fail catastrophically.

% The field structure similarly affects conditioning through $\kappa(\mathbf{J})$. Uniform or near-uniform field regions pose particular challenges as $\det(\mathbf{J}) \to 0$, requiring alternative control strategies when detected.

% \subsection{Robustness and Sensitivity Implications}

% The estimation framework demonstrates several favorable characteristics for practical deployment:

% \begin{itemize}
% \item \textbf{Graceful degradation}: Formation perturbations increase $\kappa(\mathbf{A})$ gradually, maintaining functionality with reduced accuracy rather than catastrophic failure.
% \item \textbf{Noise averaging}: Additional robots beyond the minimum three provide noise reduction through least-squares averaging as well as redundancy for noise rejection, anomaly detection and correction.
% \end{itemize}

% Key limitations require careful consideration:

% \begin{itemize}
% \item \textbf{Systematic bias}: Common-mode sensor errors translate directly to estimation bias without attenuation.
% \item \textbf{Singular field regions}: Near uniform fields where $\det(\mathbf{J}) \approx 0$, estimation becomes unreliable, necessitating fallback strategies.
% \item \textbf{Formation size constraints}: The $\mathcal{O}(s^2)$ error scaling demands compact formations in highly nonlinear fields.
% \end{itemize}





\subsection{Control Primitives}
\label{sec:control_primitives}

We develop two control primitives that utilize the estimated critical point $\mathbf{p}^*$.

\subsubsection{Critical Point Attraction}
To navigate the cluster center $\mathbf{p}_c$ to the critical point, we apply a proportional control law:
\begin{equation}
\dot{\mathbf{p}}_c = k(\mathbf{p}^* - \mathbf{p}_c)
\end{equation}
where $k > 0$ is the proportional gain. This linear error dynamic guarantees global exponential convergence to the estimated critical point with time constant $\tau = 1/k$.

\subsubsection{Orbital Navigation}
To maintain a circular orbit of radius $r_d$ around $\mathbf{p}^*$, we define the radial vector $\mathbf{r} = \mathbf{p}^* - \mathbf{p}_c$ with magnitude $r = \|\mathbf{r}\|$. The unit radial vector is $\hat{\mathbf{r}} = \mathbf{r}/r$, and the tangent vector $\hat{\mathbf{r}}_\perp$ is obtained via a $-90^\circ$ rotation. The orbital velocity command is:

\begin{equation}
\dot{\mathbf{p}}_c = k_t \hat{\mathbf{r}}_\perp - k_r(r - r_d)\hat{\mathbf{r}}
\end{equation}

Here, $k_t$ determines the orbital speed (where $\text{sgn}(k_t)$ dictates direction) and $k_r > 0$ drives radial convergence.
This control law decouples radial and angular dynamics. The radial error $e_r = r - r_d$ follows $\dot{e}_r = -k_r e_r$, ensuring exponential convergence to $r_d$. As $r \to r_d$, the angular velocity stabilizes to $\omega = k_t/r_d$, producing steady circular motion.

\subsection{Estimation Accuracy and Robustness}
\label{sec:estimation_accuracy}

The accuracy of the critical point estimation is governed by the sensitivity relationship:
\begin{equation}
\label{eq:sensitivity}
\|\Delta\mathbf{p}^*\| \sim \kappa(\mathbf{A}) \cdot \kappa(\mathbf{J}) \cdot \|\boldsymbol{\eta}\|
\end{equation}

This relationship captures a two-stage error propagation chain. Sensor noise $\boldsymbol{\eta}$ in the field measurements is first amplified by the condition number of the formation geometry matrix $\kappa(\mathbf{A})$ when solving for the linear fit coefficients (Eq.~4), and then further amplified by the condition number of the estimated Jacobian $\kappa(\mathbf{J})$ when solving for the critical point location (Eq.~9). The following subsections examine each factor.

\subsubsection{Formation Geometry ($\kappa(\mathbf{A})$)}
The formation geometry impacts estimation accuracy through its condition number $\kappa(\mathbf{A})$. An equilateral triangle minimizes this value ($\kappa \approx 2.45$) among three-robot configurations, making it the optimal minimal structure. A square formation further improves conditioning ($\kappa \approx 2.16$) due to better directional isotropy. However, degenerate configurations (e.g., collinear robots) cause $\kappa(\mathbf{A}) \to \infty$, leading to estimation failure. The geometry of the robots can be chosen to optimize sampling and condition number or for other mission criteria.

\subsubsection{Field Topology ($\kappa(\mathbf{J})$)}
The condition number $\kappa(\mathbf{J})$ captures how the local field structure amplifies coefficient errors into critical point location errors. Unlike $\kappa(\mathbf{A})$, this factor is an intrinsic property of the field and is not under the designer's control. The condition number grows when the Jacobian eigenvalues have very different magnitudes, as occurs near highly asymmetric nodes or saddle points with disparate stable and unstable manifold strengths. In the limiting case, $\kappa(\mathbf{J}) \to \infty$ as $\det(\mathbf{J}) \to 0$, corresponding to non-hyperbolic critical points or regions of near-uniform flow where no isolated critical point exists. These singular cases are detectable through eigenvalue analysis of the estimated Jacobian and require fallback strategies. For the six canonical fields used in our experiments (Table~\ref{tab:field_summary}), the symmetric eigenvalue structures yield $\kappa(\mathbf{J})$ in the range of 1--3, contributing to the favorable estimation accuracy observed in both simulation and hardware validation.

\subsubsection{Sensor Noise ($\boldsymbol{\eta}$)}
In practice, sensor noise is the dominant error source. The condition numbers $\kappa(\mathbf{A})$ and $\kappa(\mathbf{J})$ act as multipliers on the measurement noise $\|\boldsymbol{\eta}\|$, so reducing noise at the sensor level has a direct effect on estimation accuracy. For a given sensor, increasing the formation scale improves the effective signal-to-noise ratio: larger spatial separations produce greater field differences between robots while the per-sensor noise floor remains fixed. However, the linear plane fit (Eqs.~2--3) assumes locally linear field behavior, so formations must remain compact enough that this approximation holds. At our experimental scale ($s = 0.33$~m), the low systematic bias observed in hardware confirms that the linear assumption is satisfied and sensor noise, rather than model mismatch, limits performance.


\subsubsection{Robustness to Estimation Uncertainty}
The control primitives in Section \ref{sec:control_primitives} are robust estimation errors. Let $\boldsymbol{\epsilon} = \mathbf{p}^*_{est} - \mathbf{p}^*_{true}$ denote the error derived from Equation \ref{eq:sensitivity}. For the attraction controller, the system converges exponentially to $\mathbf{p}^*_{est}$. Similarly, the orbital controller maintains stable circular motion around $\mathbf{p}^*_{est}$. Thus, estimation errors shift the equilibrium location but preserve the global exponential stability of the control laws.

%%%% Here we go!!!!



\section{Control Architecture}

\subsection{Three Layered Controller Overview}

The multirobot system control architecture chosen is an established layered approach successfully demonstrated in prior cluster space control studies in both 2D and 3D environments[9, 10,19]. The architecture consists of three primary layers: (1) an adaptive navigation layer that implements the critical point detection and control laws from Section II (equations 8-17), generating desired cluster centroid velocity, (2) a cluster space controller that maintains formation geometry while translating cluster-level velocity commands into individual robot velocity commands, and (3) individual robot controllers that receive velocity commands and translate them into mechanical actuation of each robot's drive system. Figure 2 illustrates this hierarchical structure.


\begin{figure}[!t]
    \centering
    \includegraphics[width=0.45\textwidth]{control_architecture_2.png}
    \caption{Hierarchical control architecture with individual robot, cluster space, and adaptive navigation layers for coordinated motion[9]}
    \label{fig:control_architecture}
\end{figure}

This architecture was chosen for its modular design. Navigation Controllers can easily be interchanged, allowing for other testing paramaters to remain consistent between trials. It was also chosen as cluster control has been proven to maintain formation better than controllers such as leader follower chains and swarm optimazation techniques. It also lends itself to being adapted to different robotic platforms, while only changing the robot controllers.




\subsection{Three-Robot Cluster Space Controller}

Consider a cluster of three robots forming a triangular configuration. The shape can be parameterized using Side-Angle-Side (SAS) representation with parameters $(p, \beta, q)$, where $p$ is the distance between robots 1 and 2, $\beta$ is the interior angle at robot 2, and $q$ is the distance between robots 2 and 3. Figure 3 shows how this is visualized.

\begin{figure}[!t]
    \centering
    \includegraphics[width=0.37\textwidth]{threerobot.png}
    \caption{Three robot triangular formation parameterized by SAS variables: side lengths p, q and angle beta [20]}
    \label{fig:threerobot}
\end{figure}

Equations for determining the cluster shape parameters p, q, and r  given robot positions $(x_1, y_1), (x_2, y_2), (x_3, y_3)$ are known as the forward kinematic equation, while going from cluster frame variables to robot variables are known as the inverse kinematic equations.
These equations can be described analytically, and can be differentiated to give analytic forms for the Jacobian and Inverse Jacobian, which translate velocities from robot space to cluster space. The full equations can be found in [31].
\subsection{Four-Robot Cluster Space Controller}

For scaling beyond the minimum 3 robot configuration needed for implementation, a four robot cluster controller was developed. A four-robot cluster forming a quadrilateral can be parameterized using diagonal-based representation with eight parameters:
\begin{itemize}
\item $d_1, d_2$: lengths of diagonals connecting robots (1,3) and (2,4) respectively
\item $r_1, r_2$: intersection ratios indicating where diagonals cross (0 = start, 1 = end)
\item $\phi$: angle between diagonals
\item $(x_c, y_c)$: centroid position
\item $\theta_c$: overall orientation (angle of diagonal $d_1$)
\end{itemize}
A visualization can be seen in Figure 4.

\begin{figure}[!t]
    \centering
    \includegraphics[width=0.37\textwidth]{robot4fig2.png}
    \caption{Four robot quadrilateral formation with diagonal parameterization: lengths d1, d2, ratios r1, r2, and $\phi$}
    \label{fig:fourrobot}
\end{figure}

A full description of the forward kinematics, inverse kinematics, and the inverse Jacobian can be found in Appendix B.

\section{Control Performance in Simulation}

\noindent The python simulation implements a three-layer control architecture discussed in the previous section. The robots in the simulations are modelled as the robots used for hardware verification. Each robot's internal velocity controller exhibits first-order response:
\begin{equation}
v(t) = v_{\text{cmd}}(1 - e^{-t/\tau})
\end{equation}
where
$\tau$ is the actuator time constant. Each robot is modelled as a single omnidirectional point mass.

The continuous actuator dynamics follow:
\begin{equation}
\dot{\mathbf{v}}(t) = -\frac{1}{\tau}\mathbf{v}(t) + \frac{1}{\tau}\mathbf{v}_{\text{des}}(t)
\label{eq:continuous}
\end{equation}
For digital implementation at control period $\Delta t$, discretization yields:

\begin{equation}
\mathbf{v}[k+1] = \alpha \mathbf{v}[k] + (1-\alpha)\mathbf{v}_{\text{des}}[k]
\label{eq:momentum}
\end{equation}
where the momentum coefficient
$\alpha = e^{-\Delta t/\tau}$.

Combining with proportional formation control $\mathbf{v}_{\text{des}}(t) = k_v(\mathbf{x}_{\text{des}}(t) - \mathbf{x}(t))$ produces the final control law:

\begin{equation}
\mathbf{v}_{\text{cmd}}[k+1] = \alpha \mathbf{v}_{\text{cmd}}[k] + (1-\alpha)k_v(\mathbf{x}_{\text{des}}[k] - \mathbf{x}[k])
\label{eq:final}
\end{equation}
The parameter $\alpha$ characterizes actuator response ($\alpha \to 1$ for sluggish actuators, $\alpha \to 0$ for fast response), while $k_v$ controls the robot-level position error gain. Note that $k_v$ operates at the individual robot control layer, distinct from the cluster-level gain $k$ in the stability analysis (Section II.C), which governs convergence to critical points. For the simulated 10 Hz control loop with $\tau = 0.3$, $\alpha \approx 0.717$. The robots also have a limit of 0.3~m/s max, and a minimum speed of 0.05~m/s needed to overcome stiction.

Each robot is given velocity commands based on the cluster command given by the control primitive, and translated by the inverse Jacobian.

To validate the control primitives across diverse field topologies, we tested convergence to critical points in six canonical vector fields using noise-free simulations. Figure \ref{fig:Allfieldstocenter} demonstrates successful navigation from random starting positions in each field. Equations used to model each of the 6 fields tested can be found in Appendix A.

\begin{figure}[!t]
    \centering
    \includegraphics[width=0.37\textwidth]{figure_5.png}
    \caption{Successful navigation to critical points from 100 random starting positions across six canonical vector field types}
    \label{fig:Allfieldstocenter}
\end{figure}

Table~\ref{tab:simulation_estimation_results} summarizes the simulation performance across all six canonical field types. The simulation results achieved 100 percent success rate across 100 trials in convering to the critical point  across all six topological field types.  Systematic bias, remained below 0.002 m across all configurations. Precision, ranged from 0.013 m to 0.017 m, demonstrating consistent performance across field types.

For the three and four robot configurations, the dimensions were chosen as 0.33 m side length for the equilateral triangle, 0.3 m diagonal for the square robot cluster configuration. This was done to keep the area covered by the cluster of robots similar across all formations, thereby minimizing the effect of patch size on sampling performance.

% OLD TABLE WITH 4 DECIMAL PLACES - COMMENTED OUT
%\begin{table}[!t]
%\centering
%\caption{Simulation results across six field types. All methods achieved 100\% success over 100 trials.}
%\label{tab:simulation_estimation_results}
%\begin{tabular}{llccc}
%\hline
%\textbf{Environment} & \textbf{Cluster Size} & \textbf{Avg. Rest Point} & \textbf{Systematic} & \textbf{Precision} \\
% & & \textbf{(x, y) m} & \textbf{Bias (m)} & \textbf{$\sigma$ (m)} \\
%\hline
%\multicolumn{5}{l}{\textit{Stable Node (Sink)}} \\
%\hline
% & 3-Robot & (-0.0001, -0.0004) & 0.0004 & 0.0152 \\
% & 4-Robot & (-0.0015, -0.0011) & 0.0019 & 0.0145 \\
%\hline
%\multicolumn{5}{l}{\textit{Unstable Node (Source)}} \\
%\hline
% & 3-Robot & (0.0007, -0.0009) & 0.0011 & 0.0173 \\
% & 4-Robot & (0.0007, -0.0024) & 0.0025 & 0.0140 \\
%\hline
%\multicolumn{5}{l}{\textit{Center (Vortex)}} \\
%\hline
% & 3-Robot & (0.0000, -0.0001) & 0.0001 & 0.0171 \\
% & 4-Robot & (-0.0007, -0.0001) & 0.0007 & 0.0133 \\
%\hline
%\multicolumn{5}{l}{\textit{Stable Spiral (Sinking Vortex)}} \\
%\hline
% & 3-Robot & (-0.0005, 0.0004) & 0.0006 & 0.0128 \\
% & 4-Robot & (-0.0008, 0.0011) & 0.0014 & 0.0157 \\
%\hline
%\multicolumn{5}{l}{\textit{Unstable Spiral (Spewing Vortex)}} \\
%\hline
% & 3-Robot & (-0.0009, -0.0006) & 0.0011 & 0.0157 \\
% & 4-Robot & (-0.0005, 0.0011) & 0.0012 & 0.0133 \\
%\hline
%\multicolumn{5}{l}{\textit{Saddle Point}} \\
%\hline
% & 3-Robot & (-0.0004, -0.0002) & 0.0005 & 0.0144 \\
% & 4-Robot & (-0.0007, 0.0007) & 0.0009 & 0.0164 \\
%\hline
%\end{tabular}
%\end{table}

% UPDATED TABLE WITH 3 DECIMAL PLACES
\begin{table}[!t]
\footnotesize
\centering
\caption{Simulation results across six field types. All methods achieved 100\% success over 100 trials.}
\label{tab:simulation_estimation_results}
\setlength{\tabcolsep}{4pt}
\begin{tabular}{llccc}
\hline
\textbf{Environment} & \textbf{Cluster Size} & \textbf{Avg. Rest Point} & \textbf{Systematic} & \textbf{Precision} \\
 & & \textbf{(x, y) m} & \textbf{Bias (m)} & \textbf{$\sigma$ (m)} \\
\hline
\multicolumn{5}{l}{\textit{Stable Node (Sink)}} \\
\hline
 & 3-Robot & (0.000, 0.000) & 0.000 & 0.015 \\
 & 4-Robot & (-0.002, -0.001) & 0.002 & 0.014 \\
\hline
\multicolumn{5}{l}{\textit{Unstable Node (Source)}} \\
\hline
 & 3-Robot & (0.001, -0.001) & 0.001 & 0.017 \\
 & 4-Robot & (0.001, -0.002) & 0.002 & 0.014 \\
\hline
\multicolumn{5}{l}{\textit{Center (Vortex)}} \\
\hline
 & 3-Robot & (0.000, 0.000) & 0.000 & 0.017 \\
 & 4-Robot & (-0.001, 0.000) & 0.001 & 0.013 \\
\hline
\multicolumn{5}{l}{\textit{Stable Spiral (Sinking Vortex)}} \\
\hline
 & 3-Robot & (-0.001, 0.000) & 0.001 & 0.013 \\
 & 4-Robot & (-0.001, 0.001) & 0.001 & 0.016 \\
\hline
\multicolumn{5}{l}{\textit{Unstable Spiral (Spewing Vortex)}} \\
\hline
 & 3-Robot & (-0.001, -0.001) & 0.001 & 0.016 \\
 & 4-Robot & (0.000, 0.001) & 0.001 & 0.013 \\
\hline
\multicolumn{5}{l}{\textit{Saddle Point}} \\
\hline
 & 3-Robot & (0.000, 0.000) & 0.000 & 0.014 \\
 & 4-Robot & (-0.001, 0.001) & 0.001 & 0.016 \\
\hline
\end{tabular}
\end{table}


In terms of estimation accuracy, systematic bias remained below 0.002 m across all configurations, demonstrating excellent convergence to true critical points. Precision values ranged from 0.013 m to 0.017 m for the 3-robot configuration and 0.013 m to 0.016 m for the 4-robot configuration. The configurations showed comparable performance, with no single configuration consistently outperforming the other across all field types.

Orbital control performance was evaluated in simulation across all six canonical field types at commanded radii from 0.01 m to 0.70 m. Table~\ref{tab:simulation_orbital_results} shows that both 3-robot and 4-robot configurations achieve identical tracking performance with near-zero standard deviations, confirming the theoretical control law behavior in noise-free conditions.

\begin{table}[!t]
    \centering
    \caption{Simulated orbital control across all six field types. Results identical in noise-free conditions.}
    \label{tab:simulation_orbital_results}
    \begin{tabular}{ccc}
    \hline
    \textbf{Commanded} & \textbf{3-Robot} & \textbf{4-Robot Square} \\
    \textbf{Radius (m)} & \textbf{Avg. Radius} & \textbf{Avg. Radius} \\
     & \textbf{($\pm$ std)} & \textbf{($\pm$ std)} \\
    \hline
    0.01 & 0.156 $\pm$ 3.00e-06 & 0.156 $\pm$ 3.00e-06 \\
    0.10 & 0.209 $\pm$ 7.00e-06 & 0.209 $\pm$ 7.00e-06 \\
    0.20 & 0.281 $\pm$ 1.20e-05 & 0.281 $\pm$ 1.20e-05 \\
    0.30 & 0.363 $\pm$ 1.40e-05 & 0.363 $\pm$ 1.40e-05 \\
    0.40 & 0.450 $\pm$ 1.50e-05 & 0.450 $\pm$ 1.50e-05 \\
    0.50 & 0.542 $\pm$ 1.50e-05 & 0.542 $\pm$ 1.50e-05 \\
    0.60 & 0.636 $\pm$ 1.40e-05 & 0.636 $\pm$ 1.40e-05 \\
    0.70 & 0.731 $\pm$ 1.30e-05 & 0.731 $\pm$ 1.30e-05 \\
    \hline
    \end{tabular}
    \end{table}


\section{Hardware Validation}


\subsection{Hardware Test Bed}

The testbed used for validation is fully described in [5]. Figure 6 shows the complete experimental testbed integrating the robots, printed field map, and motion capture system. The testbed's modular design enables systematic experimentation: vector fields can be easily swapped by changing printed maps, robot team size can be varied from three to four units, and different control primitives can be evaluated using the same hardware platform and downstream control infrastructure. This consistency across experiments ensures that performance differences reflect algorithmic rather than implementation variations.

\begin{figure}[!t]
    \centering
    \includegraphics[width=0.37\textwidth]{testbed_with_four.png}
    \caption{Complete experimental testbed with 1.6 m × 1.6 m printed vector field map, three Decabot rovers, and OptiTrack system}
    \label{fig:fullbed}
    \end{figure}


The experimental platform consists of three omnidirectional Decabot rovers operating within a 1.6$\times$1.6 m workspace. Each Decabot features four omnidirectional wheels providing holonomic motion capability and a downward-facing TCS34725 RGB color sensor enclosed in a light-shrouded housing for consistent readings.The robots' RGB sensors are calibrated to map raw color measurements to vector field values using neural network models trained on extensive calibration data, achieving mean angular error below 7$^\circ$ and magnitude error under 10\% across the operational range.

\begin{figure}[!t]
    \centering
    \includegraphics[width=0.37\textwidth]{color_sensor_hsv_predictions.png}
    \caption{Neural network predictions for HSV field encoding showing 7 degree hue and 5\% saturation accuracy}
    \label{fig:HuePredictions}
\end{figure}

Full details of the robot calibration process can be found in [31].

An eight-camera OptiTrack motion capture system tracks robot positions with sub-centimeter accuracy, providing ground truth localization independent of the robots' onboard sensing.


The control architecture is implemented in MATLAB/Simulink running on a central computer that communicates with robots via WiFi at 10 Hz. At each control cycle, the system receives robot positions from OptiTrack and RGB sensor readings from each Decabot. The cluster space controller computes formation-maintaining velocity commands based on the selected control primitive and transmits individual robot velocity commands back to the Decabots.






Vector fields are physically realized as printed floor maps using HSV color encoding, where hue (0-2$\pi$) represents flow direction and saturation (0-1) encodes magnitude. After creating vectors at all grid points, magnitudes were scaled to the 0.3-1.0 range for printing to ensure reliable sensor readings and avoid robot stiction at low velocities. This printed range is remapped to 0-1 in software, with saturation 1.0 corresponding to the maximum robot velocity of 0.3 m/s. This encoding provides field representation that is minimally affected by lighting variations, as the hue channel and saturation channel remain stable under different illumination conditions.

This method of representing vector fields in printed floor maps has imperfections in both printing and sensing. Figure \ref{fig:measurement_error} shows us the analytical fields of both the saddle and vortex field, as well as a reconstrunction of the field from sensed values. The bottom row of plots shows us the reconstruction error.


\begin{figure}[!t]
    \centering
    \includegraphics[width=0.45\textwidth]{measurement_error_comparison.png}
    \caption{Saddle Plot Desired Vs Actual}
    \label{fig:measurement_error}
\end{figure}


\subsection{Validation Test Plan}

We conducted extensive hardware validation focusing on two representative field types that encompass the key challenges of vector field navigation: rotational dynamics (vortex) and hyperbolic instability (saddle). The vortex field tests pure rotational control dynamics and provides continuity with previous studies in our research group. The saddle field's hyperbolic structure provides the most stringent test of our estimation accuracy, as small errors can cause robots to diverge along unstable manifolds.

For each field, we selected 8 starting locations around the perimeter and ran the center attraction primitive 10 times at each location with a 3-robot triangular formation, yielding 80 trials per field. An additional 80 vortex trials were conducted with a 4-robot square formation to demonstrate generalization across cluster sizes. Due to file management errors, 3 saddle trials and 1 vortex trial were lost, resulting in 236 convergence trials. Combined with 22 orbital control experiments at varying radii, this produced 258 hardware trials overall.

Cluster formation dimensions were chosen to maintain similar coverage area across configurations, thereby minimizing the effect of patch size on sampling performance. The equilateral triangle used a 0.33 m side length and the square used a 0.3 m diagonal. These configurations match those used in simulation.




\subsection{Convergence to Critical Points on Vortex and Saddle Fields}



Figures 9 shows how repeatable the paths traveled by the robots are, which is to be expected as sensor readings are deterministic.
The solid lines are the average path travelled, while the blurred lines show the standard deviation of the travelled path.

\begin{figure}[!t]
    \centering
    \includegraphics[width=0.37\textwidth]{allrunsadjusted.png}
    \caption{Convergence trajectories from eight starting positions showing repeatable paths to vortex center with three robots}
    \label{fig:vortextomiddle}
\end{figure}



Figure 10 shows the average distance from the centroid against time over 10 runs in the vortex field. This performance confirms the predicted exponential stability for this control primitive.

\begin{figure}[!t]
    \centering
    \includegraphics[width=0.37\textwidth]{convergence.png}
    \caption{Distance to center over time for 10 runs with average trajectory and standard deviation bounds shown. The plotted mean distance $\mathbb{E}[|r|]$ combines systematic bias and random scatter (precision $\sigma$) from Table~\ref{tab:estimation_results} as $\sqrt{\text{bias}^2 + \sigma^2}$.}
    \label{fig:Tomiddlewithtime}
\end{figure}

Table~\ref{tab:estimation_results} summarizes hardware experimental results for critical point estimation. Hardware tests demonstrated successful convergence to critical points. The 3-robot configuration achieved systematic biases of 0.005--0.012 m with precision of 0.009--0.014 m. The 4-robot configuration showed a bias-variance tradeoff: systematic bias improved to 0.010 m (17\% reduction from 3-robot vortex case) but precision degraded to 0.018 m. Figure~\ref{fig:Tomiddlewithtime} displays mean distance from origin over time, which combines these bias and precision metrics. Consequently, the saddle configuration (low bias, high precision) exhibits similar mean distance to the 3-robot vortex (moderate bias, low precision) despite differing performance characteristics, while the 4-robot configuration's high precision dominates its mean distance.

% OLD TABLE WITH 4 DECIMAL PLACES - COMMENTED OUT
%\begin{table}[!t]
%    \centering
%    \caption{Hardware experimental results. Saddle: 78 runs (3-Robot). Vortex: 79 runs (3-Robot), 80 runs (4-Robot). Bias measures mean offset from origin; precision ($\sigma$) measures repeatability.}
%    \label{tab:estimation_results}
%    \begin{tabular}{llccc}
%    \hline
%    \textbf{Environment} & \textbf{Cluster} & \textbf{Avg. Rest Point} & \textbf{Systematic} & \textbf{Precision} \\
%     & \textbf{Size} & \textbf{(x, y) m} & \textbf{Bias (m)} & \textbf{$\sigma$ (m)} \\
%    \hline
%    \multicolumn{5}{l}{\textit{Saddle Point}} \\
%    \hline
%     & 3-Robot & (0.0073, -0.0016) & 0.0075 & 0.0238 \\
%    \hline
%    \multicolumn{5}{l}{\textit{Center (Vortex)}} \\
%    \hline
%     & 3-Robot & (-0.0060, 0.0199) & 0.0208 & 0.0143 \\
%     & 4-Robot & (0.0149, -0.0061) & 0.0161 & 0.0307 \\
%    \hline
%    \end{tabular}
%    \end{table}

% UPDATED TABLE WITH 3 DECIMAL PLACES
\begin{table}[!t]
    \footnotesize
    \centering
    \caption{Hardware experimental results. Saddle: 78 runs (3-Robot). Vortex: 79 runs (3-Robot), 80 runs (4-Robot). Bias measures mean offset from origin; precision ($\sigma$) measures repeatability.}
    \label{tab:estimation_results}
    \setlength{\tabcolsep}{4pt}
    \begin{tabular}{llccc}
    \hline
    \textbf{Environment} & \textbf{Cluster} & \textbf{Avg. Rest Point} & \textbf{Systematic} & \textbf{Precision} \\
     & \textbf{Size} & \textbf{(x, y) m} & \textbf{Bias (m)} & \textbf{$\sigma$ (m)} \\
    \hline
    \multicolumn{5}{l}{\textit{Saddle Point}} \\
    \hline
     & 3-Robot & (0.004, -0.001) & 0.005 & 0.014 \\
    \hline
    \multicolumn{5}{l}{\textit{Center (Vortex)}} \\
    \hline
     & 3-Robot & (-0.004, 0.012) & 0.012 & 0.009 \\
     & 4-Robot & (0.009, -0.004) & 0.010 & 0.018 \\
    \hline
    \end{tabular}
    \end{table}

    The hardware experimental results validate the effectiveness of multirobot critical point estimation on physical robots. Both configurations successfully converged to critical points across all hardware trials. Similar systematic bias is seen between 3 and 4 robot configurations, with 4 robot configuration having less precision.

    To confirm this finding, the reconstructed field was loaded into the simulation program, and 1000 random starting points were chosen for the cluster. These starting locations were used for 3 and 4 robot clusters. Simulation results in noisy environments reflect hardware results in the noisy environment. In both cases, there was 100\% success rate. The three robot configuration had a Systematic Bias of  0.028 m and precision 0.012 m. The four robot configuration in simulation has a systematic bias of 0.028 m also, but a precision of 0.023

\subsection{Orbital Control Performance Analysis}

Figure 11 shows the actual path travelled around the critical point compared to the predicted path.

\begin{figure}[!t]
    \centering
    \includegraphics[width=0.37\textwidth]{orbit040_analytical_vs_real.png}
    \caption{Saddle Plot Desired Vs Actual}
    \label{fig:perfect and real saddle orbit}
\end{figure}

Table \ref{tab:orbital_radius_comparison} summarizes orbital control tests across all control primitives at radii from 0.01 m to 0.7 m.


\begin{table}[!t]
    \centering
    \caption{Hardware orbital control results (meters). Missing entries (--) indicate radius not tested.}
    \label{tab:orbital_radius_comparison}
    \begin{tabular}{cccc}
    \hline
    \textbf{Commanded} & \textbf{Vortex} & \textbf{Vortex} & \textbf{Saddle} \\
    \textbf{Radius (m)} & \textbf{3-Robot} & \textbf{4-Robot} & \textbf{3-Robot} \\
     & \textbf{($\pm$ std)} & \textbf{($\pm$ std)} & \textbf{($\pm$ std)} \\
    \hline
    0.01 & 0.256 $\pm$ 0.053 & 0.258 $\pm$ 0.047 & 0.124 $\pm$ 0.011 \\
    0.10 & 0.290 $\pm$ 0.012 & 0.314 $\pm$ 0.024 & 0.172 $\pm$ 0.018 \\
    0.20 & 0.343 $\pm$ 0.019 & 0.335 $\pm$ 0.027 & 0.258 $\pm$ 0.037 \\
    0.30 & 0.387 $\pm$ 0.015 & 0.401 $\pm$ 0.040 & 0.389 $\pm$ 0.059 \\
    0.40 & 0.432 $\pm$ 0.020 & 0.448 $\pm$ 0.044 & 0.490 $\pm$ 0.079 \\
    0.50 & 0.465 $\pm$ 0.023 & 0.462 $\pm$ 0.042 & 0.560 $\pm$ 0.095 \\
    0.60 & 0.528 $\pm$ 0.038 & 0.487 $\pm$ 0.049 & 0.637 $\pm$ 0.100 \\
    0.70 & 0.578 $\pm$ 0.038 & -- & -- \\
    \hline
    \end{tabular}
    \end{table}




The results reveal critical performance characteristics across robot configurations. At small orbital radii (0.010 m), both 3-Robot and 4-Robot Square configurations exhibit radius errors exceeding 0.26 m in hardware experiments.

Performance improves as orbital radius increases, with errors decreasing through 0.3-0.4 m where the 3-Robot configuration achieves errors as low as 0.037 m. At larger radii (0.6-0.7 m), radius errors become negative for both 3-Robot and 4-Robot Square configurations, indicating orbits smaller than commanded.




\section{Discussion}

\subsection{Comparison with Other Primitives}

We evaluate orbital control performance by comparing our primitive against two other control primitives in the vortex field. Figure 12 illustrates the fundamental difference between passive field-following and active orbital control. The Vector Sum approach exhibits outward spiral drift due to passive field-following behavior, as confirmed in hardware in [31].

The vector-to-scalar primitive, proposed in [25] was modified to move tangentially to the direction of steepest descent. However, this primitive showed similar drift when tested in simulation. This is because there was only directional guidance, and did not know that momentum drift outward was occuring.

In contrast, our orbital control primitive maintains stable circular trajectories at the desired radius by decoupling velocity commands from field magnitude through critical point estimation and proportional radial feedback.


\begin{figure}[!t]
    \centering
    \includegraphics[width=0.45\textwidth]{orbit_demonstration.png}
    \caption{Vector-sum control primitive Vs Critical Point Detection and control}
    \label{fig:primitive_comparison}
\end{figure}

\subsection{Error Analysis}

The RGB color sensors exhibit RMSE of 0.02 radians (7.2°) in hue and 0.05 (5\%) in saturation measurements which propagate through the critical point estimation process as described the sensitivity analysis equation.

For our equilateral formation ($\kappa(\mathbf{A}) \approx 2.45$) operating in typical vector fields ($\kappa(\mathbf{J}) \approx 1$--$3$), the sensor noise translates to measurement errors $\|\boldsymbol{\eta}\| \approx 0.015$ m/s in the $(u,v)$ velocity space when accounting for the 0.3 m/s maximum robot velocity. Through the plane-fitting algorithm, this produces critical point estimation errors of approximately 0.010--0.015 m. Simulation with added sensor noise yields consistent error magnitudes, indicating that sensor noise accounts for the majority of the simulation-reality gap.

The 4-robot hardware configuration exhibited lower systematic bias (0.0096 m vs 0.0124 m for 3-robot vortex) but higher precision scatter (0.0183 m vs 0.0085 m). This bias-variance tradeoff reflects the square formation's improved geometric conditioning ($\kappa = 2.16$ vs $2.45$) reducing systematic errors, while the overdetermined least-squares solution amplifies measurement noise from four independent sources. Simulation results with noise-free measurements show improved performance with four robots (Table~\ref{tab:simulation_estimation_results}), confirming that the precision degradation arises from real-world sensor noise propagation rather than algorithmic limitations.

sensor noise ($\boldsymbol{\eta}$) becomes the dominant error source. This creates a bias-variance tradeoff when scaling from three to four robots. While the four-robot square configuration reduces systematic bias via improved geometric conditioning, the overdetermined least-squares solution integrates noise from an additional sensor, potentially degrading precision in high-noise environments compared to the exactly determined three-robot case.


\subsection{Limitations}

Several limitations constrain the current implementation and experimental validation. First, the approach is restricted to planar (2D) vector fields, though many real-world applications involve three-dimensional flows. Extension to 3D fields would require at minimum four non-coplanar robots for Jacobian estimation.

Second, the testbed validates performance only on static fields where critical points remain fixed during operation. The static printed field approach enables precise ground truth validation but does not capture dynamic field behavior.

Third, robot team sizes were limited to three and four units. While the theoretical framework supports arbitrary team sizes through the overdetermined system approach, larger formations present practical implementation challenges including communication bandwidth, computational complexity, and formation control stability.

Finally, this work assumes a sampling rate sufficiently high that the system approximates continuous control. The robustness of the approach to lower update rates and communication dropouts was not tested.

\subsection{Future Work}

The memory-free architecture developed here avoids assumptions about field stationarity, suggesting applicability to time-varying environments where field structure evolves faster than robots can accumulate spatial data. Validating this capability in dynamic fields is a priority for future work.

Hardware testbed improvements would also strengthen validation. As discussed in Section VI, realizing the theoretical noise-averaging benefits of larger teams will require more rigorous calibration procedures to eliminate systematic offsets.

Beyond critical point detection and orbital control, vector fields contain other topological features of practical interest. Separatrices, the stable and unstable manifolds emanating from saddle points, define boundaries between distinct flow regimes and could serve as natural navigation corridors or barriers. Heteroclinic orbits, which connect different critical points through their manifolds, offer energy-efficient transit routes through complex fields. The Jacobian estimation framework developed here provides the foundation for detecting such features through analysis of eigenvector alignment between neighboring critical points.


\section{Conclusion}
\noindent

This paper presents a distributed framework for multirobot navigation in vector fields containing critical points. A Jacobian-based approach uses instantaneous measurements from three robots to identify critical point locations and types without prior field knowledge. Simulation across six canonical field types confirmed convergence of these control primitives. 
In 258 hardware trials, the system achieved a 100\% success rate with position errors under 0.015 m. Orbital control experiments maintained commanded radii within 0.02 m. This architecture scales to larger formations and supports precise standoff control for environmental monitoring and plume tracking.


\appendices

\section{Vector Field Environments}
\label{sec:appendix_vector_fields}

Six vector field environments centered at the origin were used to evaluate the multi-robot navigation algorithms. Let $r = \sqrt{x^2 + y^2}$ and $\theta = \arctan2(y, x)$.

\textbf{Vortex:} Pure rotational flow.
\begin{equation}
\mathbf{v}_{\text{vortex}}(x, y) = \begin{bmatrix} -r \sin(\theta) \\ r \cos(\theta) \end{bmatrix}
\end{equation}

\textbf{Sinking Vortex:} Spiraling inward flow.
\begin{equation}
\mathbf{v}_{\text{sink-vortex}}(x, y) = 0.4 \begin{bmatrix} -r \sin(\theta) \\ r \cos(\theta) \end{bmatrix} + 0.15 \begin{bmatrix} -x \\ -y \end{bmatrix}
\end{equation}

\textbf{Sink:} Radial inward flow.
\begin{equation}
\mathbf{v}_{\text{sink}}(x, y) = \begin{bmatrix} -x \\ -y \end{bmatrix}
\end{equation}

\textbf{Source:} Radial outward flow.
\begin{equation}
\mathbf{v}_{\text{source}}(x, y) = \begin{bmatrix} x \\ y \end{bmatrix}
\end{equation}

\textbf{Spewing Vortex:} Spiraling outward flow.
\begin{equation}
\mathbf{v}_{\text{spew-vortex}}(x, y) = -0.4 \begin{bmatrix} -r \sin(\theta) \\ r \cos(\theta) \end{bmatrix} - 0.15 \begin{bmatrix} -x \\ -y \end{bmatrix}
\end{equation}

\textbf{Saddle:} Hyperbolic saddle point rotated $45^\circ$. Coordinates are first rotated by $\pi/4$, the saddle field is applied, then rotated back by $-\pi/4$:
\begin{equation}
\mathbf{v}_{\text{saddle}}(x, y) = \begin{bmatrix} -y \\ -x \end{bmatrix}
\end{equation}
This creates a saddle with separatrices along the diagonals rather than the coordinate axes.

\section{Four-Robot Cluster Controller Equations}

\subsection{Forward Kinematics}

The forward kinematics transforms corner positions to diagonal parameters:

\begin{equation}
x_c = \frac{1}{4}\sum_{i=1}^{4} x_i, \quad y_c = \frac{1}{4}\sum_{i=1}^{4} y_i
\end{equation}

\begin{equation}
d_1 = \sqrt{(x_3 - x_1)^2 + (y_3 - y_1)^2}
\end{equation}

\begin{equation}
d_2 = \sqrt{(x_4 - x_2)^2 + (y_4 - y_2)^2}
\end{equation}

\begin{equation}
\theta_c = \text{atan2}(y_3 - y_1, x_3 - x_1)
\end{equation}

\begin{equation}
\phi = \text{atan2}(y_3 - y_1, x_3 - x_1) - \text{atan2}(y_4 - y_2, x_4 - x_2)
\end{equation}

The intersection ratios $r_1$ and $r_2$ are computed by solving the line intersection:

\begin{equation}
r_1 = \frac{(x_2 - x_1)(y_2 - y_4) + (y_2 - y_1)(x_4 - x_2)}{(x_3 - x_1)(y_2 - y_4) + (y_3 - y_1)(x_4 - x_2)}
\end{equation}

\begin{equation}
r_2 = \frac{(x_3 - x_1)(y_2 - y_1) - (y_3 - y_1)(x_2 - x_1)}{(x_3 - x_1)(y_2 - y_4) + (y_3 - y_1)(x_4 - x_2)}
\end{equation}

\subsection{Inverse Kinematics}

Given the diagonal parameters, the robot positions are recovered as follows. First, compute the intersection point:

\begin{equation}
\begin{aligned}
x_i &= x_c + \frac{1}{4}\left[d_1(2r_1 - 1)\cos(\theta_c) + d_2(2r_2 - 1)\cos(\theta_c - \phi)\right] \\
y_i &= y_c + \frac{1}{4}\left[d_1(2r_1 - 1)\sin(\theta_c) + d_2(2r_2 - 1)\sin(\theta_c - \phi)\right]
\end{aligned}
\end{equation}

Then compute corner positions:

\begin{equation}
\begin{aligned}
x_1 &= x_i - r_1 d_1 \cos(\theta_c) \\
y_1 &= y_i - r_1 d_1 \sin(\theta_c)
\end{aligned}
\end{equation}

\begin{equation}
\begin{aligned}
x_2 &= x_i - r_2 d_2 \cos(\theta_c - \phi) \\
y_2 &= y_i - r_2 d_2 \sin(\theta_c - \phi)
\end{aligned}
\end{equation}

\begin{equation}
\begin{aligned}
x_3 &= x_i + (1 - r_1) d_1 \cos(\theta_c) \\
y_3 &= y_i + (1 - r_1) d_1 \sin(\theta_c)
\end{aligned}
\end{equation}

\begin{equation}
\begin{aligned}
x_4 &= x_i + (1 - r_2) d_2 \cos(\theta_c - \phi) \\
y_4 &= y_i + (1 - r_2) d_2 \sin(\theta_c - \phi)
\end{aligned}
\end{equation}

\subsection{Inverse Jacobian}

Following the same structure as the three-robot inverse Jacobian, the error propagation from shape parameters to robot positions extends to four robots:

\begin{equation}
\begin{bmatrix}
\delta x_1 \\ \delta y_1 \\ \delta x_2 \\ \delta y_2 \\ \delta x_3 \\ \delta y_3 \\ \delta x_4 \\ \delta y_4
\end{bmatrix} =
\mathbf{J}^{-1}
\begin{bmatrix}
\delta x_c \\ \delta y_c \\ \delta \theta_c \\ \delta d_1 \\ \delta d_2 \\ \delta r_1 \\ \delta r_2 \\ \delta \phi
\end{bmatrix}
\end{equation}

where $\mathbf{J}^{-1}$ is the $8 \times 8$ inverse Jacobian matrix containing partial derivatives $\frac{\partial x_i}{\partial \xi_j}$ for position coordinates $x_i$ and shape parameters $\xi_j$, analogous to the three-robot formulation but extended to accommodate the additional robot and diagonal parameterization.


\begin{thebibliography}{99}

\bibitem{1} A. S. Matveev et al., "Nongradient-based tracking of environmental field isolines in a changing medium by dubins-vehicle type robots," \textit{IEEE Transactions on Automatic Control}, vol. 69, no. 2, pp. 952--967, 2023.

\bibitem{2} M. Hoy, A. S. Matveev, and A. V. Savkin, "Algorithms for collision-free navigation of mobile robots in complex cluttered environments: a survey," \textit{Robotica}, vol. 33, no. 3, pp. 463--497, 2015.

\bibitem{3} X. Chen and J. Huang, "Odor source localization algorithms on mobile robots: A review and future outlook," \textit{Robotics and Autonomous Systems}, vol. 112, pp. 123--136, 2019.

\bibitem{4} J. Cortés and M. Egerstedt, "Coordinated control of multi-robot systems: A survey," \textit{SICE Journal of Control, Measurement, and System Integration}, vol. 10, no. 6, pp. 495--503, 2017.

\bibitem{5} H. Chen et al., "Localization for multirobot formations in indoor environment," \textit{IEEE/ASME Transactions on Mechatronics}, vol. 15, no. 4, pp. 561--574, 2009.

\bibitem{6} H. Xie et al., "Autonomous multirobot navigation and cooperative mapping in partially unknown environments," \textit{IEEE Transactions on Instrumentation and Measurement}, vol. 72, pp. 1--12, 2023.

\bibitem{7} J. Zhang et al., "Multi-AUV adaptive path planning and cooperative sampling for ocean scalar field estimation," \textit{IEEE Transactions on Instrumentation and Measurement}, vol. 71, pp. 1--14, 2022.

\bibitem{8} H. J. S. Feder, J. J. Leonard, and C. M. Smith, "Adaptive mobile robot navigation and mapping," \textit{The International Journal of Robotics Research}, vol. 18, no. 7, pp. 650--668, 1999.

\bibitem{9} C. A. Kitts, R. T. McDonald, and M. A. Neumann, "Adaptive navigation control primitives for multirobot clusters: Extrema finding, contour following, ridge/trench following, and saddle point station keeping," \textit{IEEE Access}, vol. 6, pp. 17625--17642, 2018.

\bibitem{10} R. T. McDonald, C. A. Kitts, and M. A. Neumann, "Navigation of scalar fronts with multirobot clusters in simulation and experiment," \textit{IEEE Systems Journal}, vol. 14, no. 3, pp. 3755--3766, 2020.

\bibitem{11} P. Ögren, E. Fiorelli, and N. E. Leonard, "Cooperative control of mobile sensor networks: Adaptive gradient climbing in a distributed environment," \textit{IEEE Transactions on Automatic Control}, vol. 49, no. 8, pp. 1292--1302, 2004.

\bibitem{12} F. Pagano et al., "Distributed Multi-Robot Active-Sensing of a Diffusive Source," \textit{IEEE Robotics and Automation Letters}, 2025.

\bibitem{13} Z. Deng and J. Luo, "Fully distributed algorithms for constrained nonsmooth optimization problems of general linear multiagent systems and their application," \textit{IEEE Transactions on Automatic Control}, vol. 69, no. 2, pp. 1377--1384, 2023.

\bibitem{14} L. Briñón-Arranz, A. Renzaglia, and L. Schenato, "Multirobot symmetric formations for gradient and hessian estimation with application to source seeking," \textit{IEEE Transactions on Robotics}, vol. 35, no. 3, pp. 782--789, 2019.

\bibitem{15} Y. Deng, B. Zhu, and H. Duan, "Bioinspired bearing–based target enclosing control for unmanned aerial vehicle swarm," \textit{IEEE/ASME Transactions on Mechatronics}, 2024.

\bibitem{16} J. Wang et al., "Dynamic plume tracking by cooperative robots," \textit{IEEE/ASME Transactions on Mechatronics}, vol. 24, no. 2, pp. 609--620, 2019.

\bibitem{17} T. Adamek, C. A. Kitts, and I. Mas, "Gradient-based cluster space navigation for autonomous surface vessels," \textit{IEEE/ASME Transactions on Mechatronics}, vol. 20, no. 2, pp. 506--518, 2014.

\bibitem{18} Z. Jin et al., "Gradient-free cooperative source-seeking of quadrotor under disturbances and communication constraints," \textit{IEEE Transactions on Industrial Electronics}, 2024.

\bibitem{19} R. K. Lee et al., "Multiple UAV adaptive navigation for three-dimensional scalar fields," \textit{IEEE Access}, vol. 9, pp. 122626--122654, 2021.

\bibitem{20} J. Burgués and S. Marco, "Environmental chemical sensing using small drones: A review," \textit{Science of the total environment}, vol. 748, p. 141172, 2020.

\bibitem{21} I. V. Bychkov, A. A. Tolstikhin, and S. A. Ulyanov, "A lobster-inspired multi-robot control strategy for monitoring non-stationary concentration fields," \textit{Journal of Physics: Conference Series}, vol. 2099, no. 1, IOP Publishing, 2021.

\bibitem{22} Z. Zhou, J. Liu, and J. Yu, "A survey of underwater multi-robot systems," \textit{IEEE/CAA Journal of Automatica Sinica}, vol. 9, no. 1, pp. 1--18, 2021.

\bibitem{23} D. Panagou, "Motion planning and collision avoidance using navigation vector fields," in \textit{Proc. IEEE International Conference on Robotics and Automation}, 2014, pp. 2513--2518.

\bibitem{24} X. Wang et al., "Adaptive vector field guidance without a priori knowledge of course dynamics and wind," \textit{IEEE/ASME Transactions on Mechatronics}, vol. 27, no. 6, pp. 4597--4607, 2022.

\bibitem{25} R. Cooper et al., "Initial Study of Multirobot Adaptive Navigation for Exploring Environmental Vector Fields," in \textit{Proc. IEEE Aerospace Conference}, 2019.

\bibitem{26} V. M. Gonçalves, R. McLaughlin, and G. A. S. Pereira, "Precise landing of autonomous aerial vehicles using vector fields," \textit{IEEE Robotics and Automation Letters}, vol. 5, no. 3, pp. 4337--4344, 2020.

\bibitem{27} B. Hu et al., "Spontaneous-ordering platoon control for multirobot path navigation using guiding vector fields," \textit{IEEE Transactions on Robotics}, vol. 39, no. 4, pp. 2654--2668, 2023.

\bibitem{28} D. R. Nelson et al., "Vector field path following for miniature air vehicles," \textit{IEEE Transactions on Robotics}, vol. 23, no. 3, pp. 519--529, 2007.

\bibitem{29} A. S. Matveev and A. A. Semakova, "Distributed 3D navigation of swarms of non-holonomic UAVs for coverage of unsteady environmental boundaries," \textit{Drones}, vol. 6, no. 2, p. 33, 2022.

\bibitem{30} C. A. Kitts and I. Mas, "Cluster space specification and control of mobile multirobot systems," \textit{IEEE/ASME Transactions on Mechatronics}, vol. 14, no. 2, pp. 207--218, 2009.

\bibitem{31} C. Waight and C. Kitts, "A functional indoor testbed for multirobot adaptive navigation in vector field environments," in \textit{Proc. ASME Int. Design Engineering Technical Conf. Computers and Information in Engineering Conf.}, vol. 89251, 2025.

\bibitem{32} R. T. McDonald, M. Condino, M. A. Neumann, and C. A. Kitts, "Navigation of scalar fronts with multirobot clusters in simulation and experiment," \textit{IEEE Systems Journal}, vol. 14, no. 3, pp. 3755--3766, 2020.


\end{thebibliography}


\begin{IEEEbiography}[{\includegraphics[width=1in,height=1.25in,clip,keepaspectratio]{fig1.jpg}}]{Christopher Waight}
received his BSc Degree in Physics from Santa Clara University. He later earned his MSc in Mechanical Engineering there, and is now a PhD Candidate in Mechanical Engineering. He currently works as a Research Assistant in SCU's Robotics System Lab. He also has experience working as a Technical Program Manager in R\&D.
\end{IEEEbiography}

\begin{IEEEbiography}[{
\includegraphics[width=1in,height=1.25in,clip,keepaspectratio]{ckitts.png}}]{Christopher A. Kitts}
    (Senior Member, IEEE) earned degrees from Princeton University (B.S.E.), University of Colorado (M.P.A.), and Stanford University (M.S. and Ph.D.). He served as a Satellite Constellation Mission Controller with the U.S. Air Force and as a Computer Scientist at NASA Ames Research Center. He is currently Professor of Mechanical Engineering, Director of the Robotic Systems Laboratory, and Associate Dean at Santa Clara University. Dr. Kitts is a Fellow of ASME.
    \end{IEEEbiography}

\end{document}
